\documentclass[a4paper,11pt]{article}
\usepackage[margin=1.5cm]{geometry}
\usepackage{multicol}
\usepackage{listings}
\usepackage{xcolor}

\lstset{
  basicstyle=\ttfamily\small,
  keywordstyle=\color{blue},
  commentstyle=\color{gray}
}

\begin{document}

\begin{center}
    {\Large \textbf{Formulario C++ Básico}}\\
\end{center}

% \begin{multicols}{2}

\section*{Estructura básica}
\begin{multicols}{2}
    %lstlisting muestra código fuente 
\begin{lstlisting}  
  
#include <iostream> 
#include "iostream"
using namespace std;
int main() {
    cout << "Hola";
    return 0;
}
\end{lstlisting}

\columnbreak

La diferencia entre 
\begin{lstlisting}
#include <iostream> e #include "filename" 
\end{lstlisting}

es la forma en la que el compilador busca los archivos. 
Con \textbf{$<$ $>$} el compilador busca el archivo en una dirección estandar, 
mientras que con \textbf{$"$ $"$}, el archivo se busca primero en el directorio del proyecto.
Si no se encuentra ahí, entonces lo busca como con \textbf{$<$ $>$}.

\end{multicols}


\section*{Tipos de Datos}
\begin{center}
\begin{tabular}{| c | c | c | c |}
        \hline
        Type & Size (bytes) & Ejemplo & Descripción\\
        \hline\hline
        int & 4 bytes & int x = 10 & Números enteros \\
        \hline
        char & 1 byte & char c = 'A' & Un sólo caracter \\
        \hline
        short & 2 bytes & short s = 100 & Entero pequeño \\
        \hline
        long & 4-8 bytes & long l = 100000 & Entero grande \\
        \hline
        float & 4 bytes & float f = 10.5 & Decimal Single-precision \\
        \hline 
        double & 8 bytes & double d = 10.55 & Decimal Double-precision \\
        \hline
        bool & 1 byte & bool b = true & True/False Values \\
        \hline
        pointer & 4-8 bytes & int * p = \&v ('v' being an int type) & Guarda dirección de memoria de una variable\\
        \hline
        void & - & - & Representa la falta de valor \\
        \hline
\end{tabular}
\end{center}
    
\section*{Operaciones Aritméticas}
\begin{center}
    \begin{tabular}{| c | c | c | c |}
        \hline
        Operación & Símbolo & Ejemplo & Asignación Compuesta \\
        \hline\hline
        Suma & + & a + b & a+=b $\rightarrow$  a = a+b\\
        \hline
        Resta & - & a - b & a-=b $\rightarrow$ a = a-b\\
        \hline
        Multiplicación & * &  a * b & a*=b $\rightarrow$ a = a*b\\
        \hline
        División & / & a / b & a/=b $\rightarrow$ a = a/b\\
        \hline
        Modulo & \% & a \% b & a\%=b $\rightarrow$ a = a\%b\\
        \hline
        Potencia & NA &$<$cmath$>$ std::pow(a, 3) & NA \\
        \hline
    \end{tabular}
\end{center}

\section*{Comparadores}
\begin{center}
    \begin{tabular}{| c | c |}
        \hline
        Operador & Descripción \\
        \hline\hline
        == & Igual a \\
        \hline
        != & NO igual a \\
        \hline
        $<$ & Menor que \\
        \hline
        $>$ & Mayor que \\
        \hline
        $<$= & Menor o igual que \\
        \hline
        $>$= & Mayor o igual que \\
        \hline
    \end{tabular}
\end{center}

\section*{Operaciones lógicas}
\begin{center}
    \begin{tabular}{| c | c | c |}
    
    \end{tabular}
\end{center}

\section*{Operaciones lógicas BITWISE (Sólo variables enteras)}
\begin{center}
    \begin{tabular}{| c | c | c |}
    \hline
    Operador & Operación lógica & Ejemplo a 8 bits \\
    \hline\hline 
    \& & Bitwise AND & 00001100 (12) \& 00011001 (25) = 00001000 (8) \\
    \hline
    $|$ & Bitwise OR &  00001100 (12) $|$ 00011001 (25) = 00011101 (29) \\
    \hline
    $\wedge$ & Bitwise XOR &  00001100 (12) $\wedge$ 00011001 (25) = 00010101 (21) \\
    \hline
    $\sim$  & Bitwise Complemento  & $\sim$ 00001100 (12) = 11110011 (-13)\\
    \hline
    $<<$ & Left Shift & 00001100 (12) $<<$ 1 = 00011000 (24)\\
    \hline
    $>>$ & Right Shift &  00001100 (12) $>>$ 1 = 00001100 (6) \\
    \hline
    \end{tabular}
\end{center}

Con estos operadores también se puede hacer la Asignación Compuesta.\\

Por sí solo, C++ no permite usar estos operadores con floats o doubles, pero se puede hacer
utilizando \begin{lstlisting}
    std::bitcast (C++ 20, <bit>)
\end{lstlisting} 
para obtner la representación binaria del número real y poder usar los 
operadores de BITWISE (No es común, pero se puede).

\section*{Control}
\begin{lstlisting}
if (COND) {}
for (int i=0;i<n;i++) {}
while(COND) {}
\end{lstlisting}

\section*{Funciones}
\begin{lstlisting}
int suma(int a, int b) {
    return a + b;
}
\end{lstlisting}

\section*{Clases}
\begin{lstlisting}
class A {
public:
    int x;
    A(int v): x(v) {}
};
\end{lstlisting}

\section*{STL}
\begin{itemize}
    \item \texttt{vector<T>}
    \item \texttt{map<K,V>}
    \item \texttt{set<T>}
    \item \texttt{unordered\_map}
\end{itemize}

\section*{Compilación}
\begin{itemize}
    \item \texttt{g++ main.cpp -o main}
    \item \texttt{g++ -std=c++17 file.cpp}
\end{itemize}

% \end{multicols}
\end{document}
